\chapter{\MakeUppercase{Introdução}}

Incêndios florestais constituem um problema de grande relevância ambiental,
econômica e social. Além de afetarem a vegetação, a biodiversidade e os serviços
ecossistêmicos, esses eventos podem ameaçar comunidades, infraestruturas e
atividades produtivas. Em diversas regiões, as mudanças climáticas tendem a
intensificar condições meteorológicas favoráveis à ocorrência de incêndios
\cite{IPCC2023}. Esse cenário reforça a necessidade de estratégias de prevenção,
planejamento e resposta.

Durante uma operação de supressão, a disponibilidade limitada de equipes,
equipamentos e tempo impede, em geral, que todas as regiões ameaçadas sejam
atendidas simultaneamente. Torna-se necessário decidir onde e quando empregar
os recursos, considerando que cada intervenção pode modificar a propagação do
fogo e, consequentemente, a viabilidade e a qualidade das decisões posteriores.
Problemas dessa natureza motivam modelos de otimização e métodos computacionais
para estudar a alocação de recursos de combate
\cite{Alvelos2018,MendesAlvelos2023}. Neste trabalho, esse processo decisório é investigado por meio de
abordagens exatas e heurísticas, comparadas quanto à qualidade das soluções
e ao custo computacional. Essa comparação permite analisar sistematicamente
diferentes configurações do problema.

Neste trabalho, a paisagem é discretizada em células ou regiões e representada
por um grafo direcionado. Os vértices correspondem a essas unidades do terreno,
e os arcos indicam possíveis direções de propagação do fogo. Os pesos dos arcos expressam
tempos de deslocamento da frente de fogo e podem refletir fatores como vento,
inclinação e características da vegetação. A partir de um vértice de ignição, os
tempos de chegada do fogo são aproximados por caminhos mínimos; em instantes
predeterminados, recursos de supressão tornam-se disponíveis para alocação em
vértices ainda não atingidos (\citeauthor{Alvelos2018}, \citeyear{Alvelos2018}; \citeauthor{DelazeriRitt2024}, \citeyear{DelazeriRitt2024,DelazeriRitt2026}).

Esse modelo define o Problema de Supressão de Incêndios Florestais
(WSP, do inglês: \textit{wildfire suppression problem}), cujo objetivo
é encontrar uma alocação temporalmente viável dos recursos que minimize o
número de vértices queimados até um horizonte de planejamento. A aplicação de
um recurso a um vértice acrescenta um atraso à propagação pelos seus arcos de
saída, podendo alterar os tempos de chegada do fogo em outras partes do grafo.
Como as decisões são interdependentes e o número de alocações possíveis cresce
rapidamente com o tamanho da instância, encontrar uma solução ótima é
computacionalmente difícil. \citeonline{DelazeriRitt2026} demonstram que a versão de decisão do WSP
é NP-completa por redução polinomial do Problema dos Nós Mais Vitais
(MVNP, do inglês: \textit{most vital nodes problem}), fortemente
NP-completo segundo \citeonline{BarNoyKhullerSchieber1995}. A dificuldade
do WSP persiste mesmo quando todos os recursos de supressão são liberados
simultaneamente.

O WSP foi introduzido na literatura por meio de uma formulação de
Programação Inteira Mista (MIP, do inglês: \textit{mixed-integer programming})
\cite{Alvelos2018}, resolvida por \textit{solvers} de propósito geral.
Posteriormente, foram propostas a Busca Local Iterada (ILS, do inglês:
\textit{iterated local search}) \cite{MendesAlvelos2023}, a Decomposição de
Benders Baseada em Lógica (LBBD, do inglês: \textit{logic-based benders
decomposition}) \cite{HarrisForbesTaimre2023} e a Busca em Feixe Iterada
(IBS, do inglês: \textit{iterated beam search}) \cite{DelazeriRitt2024}.
O IBS explora uma árvore de alocações parciais, mantém um conjunto limitado
de soluções promissoras e prioriza regiões próximas à frente de avanço do
fogo. Nos experimentos de \citeonline{DelazeriRitt2024}, o método alcançou soluções
ótimas em quase todas as instâncias utilizadas na literatura, com tempo
computacional consideravelmente menor que o das abordagens comparadas.

Este trabalho considera também a Heurística Baseada em Cortes (CUTH,
do inglês: \textit{cut-based heuristic}) \cite{Delazeri2025}, uma
matheurística que combina um corte estático de vértices com uma alocação
gulosa ao longo do tempo, e a Busca Aleatória (RS, do inglês:
\textit{random search}), utilizada como linha de base de exploração
aleatória viável \cite{DelazeriRitt2026}. A comparação envolve, assim,
seis abordagens: os métodos exatos MIP e LBBD, os heurísticos ILS, IBS e RS
e a matheurística CUTH.

\citeonline{DelazeriRitt2026} apresentaram uma nova formulação de programação
inteira mista e propuseram um gerador de instâncias fundamentado em um modelo
físico de propagação. Os autores também mostraram que os \textit{benchmarks}
clássicos eram pequenos, pouco diversos e pouco desafiadores para os métodos
atuais. O novo gerador permite variar fatores como tamanho da grade, condições
ambientais e configuração dos recursos, produzindo casos capazes de revelar regimes de
desempenho antes ocultos pelas limitações da base clássica.

Dentre as abordagens consideradas, o IBS é adotado neste trabalho como método
de referência. Essa escolha decorre dos resultados de
\citeonline{DelazeriRitt2026}, que mostraram experimentalmente que esse
algoritmo obteve o melhor desempenho entre os métodos comparados. Por ser
estocástico e não provar otimalidade, o IBS dificilmente pode ser compreendido
apenas pelo valor final de suas soluções. Sua estrutura de buscas em feixe
sucessivas permite registrar as soluções devolvidas a cada iteração, o que
viabiliza a análise de suas trajetórias. O desempenho dos seis métodos é
comparado sob um mesmo protocolo, sem tratamento privilegiado do IBS.

Os resultados de \citeonline{DelazeriRitt2026} sobre os \textit{benchmarks}
clássicos evidenciam que eles cobrem apenas uma parcela restrita dos casos
produzíveis pelo modelo. À luz desse diagnóstico e dos princípios da avaliação por
espaços de instâncias \cite{MunozEtAl2018,SmithMilesMunoz2023}, identificam-se
três lacunas experimentais relacionadas: a cobertura limitada da coleção; a
caracterização parcial da estrutura dos grafos; e a dependência de indicadores
finais de qualidade e tempo. Esses indicadores são essenciais para comparar os
métodos, mas, isoladamente, não permitem associar o desempenho às propriedades
das instâncias nem reconstruir a dinâmica da busca que levou aos resultados
\cite{OchoaMalanBlum2021,SmithMilesMunoz2023}.

A Análise do Espaço de Instâncias (ISA, do inglês:
\textit{instance space analysis}) permite tratar a primeira lacuna ao
representar cada instância por suas características e relacionar sua posição no
espaço ao desempenho dos algoritmos. Assim, é possível examinar a diversidade dos
casos de teste, identificar diferentes níveis de dificuldade, delimitar regiões
de competência dos métodos e localizar porções insuficientemente cobertas
\cite{MunozEtAl2018,SmithMilesMunoz2023}. Neste trabalho, a ISA é operacionalizada
por meio do fluxo associado ao \textit{framework}
\textit{Melbourne Algorithm Test Instance Library
with Data Analytics} (MATILDA), que integra preparação dos dados, seleção de
atributos, projeção, predição de desempenho e construção de \textit{footprints}
\cite{SmithMilesMunoz2023}.

Para tratar a segunda lacuna, a caracterização pode ser enriquecida por medidas
estruturais de grafos. Grafos com
dimensões semelhantes podem apresentar distribuições distintas de graus,
centralidade e conectividade, propriedades que descrevem a organização de
caminhos e regiões estruturalmente relevantes \cite{Newman2018}. Incorporá-las à
ISA permite avaliar se a estrutura direcionada acrescenta informação aos
atributos originais do problema, sem pressupor melhora preditiva para todos os
algoritmos.

Por fim, as Redes de Trajetórias de Busca (STNs, do inglês:
\textit{search trajectory networks}) tratam a terceira lacuna ao representar
o comportamento de um algoritmo de busca como uma rede direcionada
\cite{ChaconSartoriEtAl2023,OchoaMalanBlum2021}. Enquanto a ISA indica onde uma
instância se posiciona no espaço e as medidas topológicas descrevem como seu
grafo se organiza, as STNs permitem examinar como o algoritmo se comporta durante
a busca. Seus nós correspondem a soluções representativas registradas ao longo das
execuções, e suas arestas registram transições consecutivas entre essas soluções.
Essa modelagem evidencia convergência, recorrência e sobreposição entre
execuções. Neste trabalho, a técnica é aplicada ao IBS.

\section{Pergunta de pesquisa e hipótese}\label{sec:pergunta-pesquisa}

A pesquisa é orientada pela seguinte pergunta: que relações podem ser
identificadas entre as propriedades das instâncias do WSP, o desempenho
dos algoritmos e as trajetórias de busca do IBS quando a cobertura desse
espaço é ampliada de forma orientada?
A questão articula a caracterização das instâncias pela ISA e pelas medidas
topológicas à análise das trajetórias do IBS. A cobertura é avaliada no
mapa fixado para a geração; as relações com desempenho são examinadas
na coleção ampliada e nos modelos ISA finais.

A hipótese deste trabalho é que a ampliação orientada da cobertura do
espaço de instâncias do WSP, aliada à caracterização por atributos do
problema e medidas topológicas, torna observáveis relações entre as
propriedades das instâncias, o desempenho dos algoritmos e o comportamento
de busca do IBS que a coleção original não permitia distinguir.
A hipótese é examinada pelo contraste entre as coleções, pelos modelos ISA
e pelos descritores das STNs. Esses contrastes buscam evidências descritivas
e associações; não isolam efeitos causais da geração ou da topologia.

\section{Objetivos}\label{sec:objetivos}

O objetivo geral é caracterizar, interpretar e ampliar a cobertura de uma
representação do espaço de instâncias do WSP, relacionando seus atributos
ao desempenho dos algoritmos e às trajetórias de busca do IBS.
Para alcançá-lo, estabeleceram-se os seguintes objetivos específicos:
\begin{alineas}
 \item consolidar e caracterizar a base experimental;
 \item avaliar a contribuição descritiva e preditiva interna das medidas topológicas;
 \item construir e interpretar a ISA com a MATILDA;
 \item gerar e validar instâncias para regiões internas insuficientemente cobertas do mapa de referência;
 \item comparar as seis abordagens na coleção ampliada;
 \item caracterizar as trajetórias do IBS, respeitando as limitações de comparabilidade entre os protocolos.
\end{alineas}

\section{Contribuição e organização do trabalho}\label{sec:contribuicao}

A principal contribuição consiste em integrar técnicas existentes de
caracterização de instâncias, análise de desempenho, povoamento orientado de
uma representação fixa e investigação de trajetórias em um fluxo experimental
aplicado ao WSP. A análise empírica resultante situa o desempenho dos métodos em
diferentes tipos de instância e examina o comportamento do IBS. Trata-se,
contudo, de uma investigação computacional sobre instâncias de \textit{benchmark}:
a validade das instâncias sintéticas
refere-se à consistência com o gerador e com os critérios adotados, não à
validação de cenários reais de incêndio ou à prescrição direta de decisões
operacionais.

A pesquisa foi organizada em três etapas. Primeiro, consolidaram-se 55 instâncias
do WSP, sendo 16 clássicas e 39 produzidas pelo gerador de
\citeonline{DelazeriRitt2026}, com resultados históricos de seis abordagens e STNs
do IBS nas instâncias clássicas. Em seguida, um mapa com atributos originais e topológicos forneceu alvos
em regiões internas pouco ocupadas. Um algoritmo genético inspirado em
\citeonline{SmithMilesBowly2015} orientou a geração de 19 novas instâncias
sintéticas, ampliando a coleção para 74 casos. Na etapa final,
as seis abordagens foram avaliadas na coleção ampliada sob um protocolo homogêneo,
a MATILDA foi recalculada com a matriz completa de desempenho e novas STNs do IBS
foram construídas. As STNs clássicas permaneceram como contexto descritivo devido
às diferenças de protocolo e ambiente computacional.

O trabalho segue com a fundamentação teórica, a metodologia, a apresentação
e discussão dos resultados e as conclusões, que retomam os objetivos e as
limitações da pesquisa.
